% Ahmad Bensalah - CV DevOps Engineer (French)
% Compile with: pdflatex CV_Ahmad_Bensalah_FR.tex

\documentclass[11pt,a4paper,sans]{moderncv}

% ModernCV theme
\moderncvstyle{banking}
\moderncvcolor{blue}

% Character encoding
\usepackage[utf8]{inputenc}
\usepackage[T1]{fontenc}
\usepackage[french]{babel}

% Adjust page margins
\usepackage[scale=0.85]{geometry}
\setlength{\hintscolumnwidth}{3cm}

% Personal data
\name{Ahmad}{Bensalah}
\title{Ingénieur DevOps \& Platform Engineering}
\phone[mobile]{+216~XX~XXX~XXX}
\email{ahmad.bensalah.contact@gmail.com}
\social[linkedin]{ahmad-bensalah}
\social[github]{ahmad-bensalah}
\homepage{portfolio-abs.pages.dev}
\extrainfo{Disponible : Février 2027 (PFE 6 mois)}

\begin{document}

\makecvtitle

% ================================
% PROFIL
% ================================
\section{Profil}
\small
Étudiant ingénieur en \textbf{Data Engineering} à la FSS, diplômé prévu en 2027, à la recherche d'un \textbf{PFE de 6 mois à partir de février 2027} en \textbf{DevOps, Platform Engineering ou MLOps}.

Je construis les \textbf{infrastructures sur lesquelles reposent les systèmes de données et d'IA}. Spécialisé en \textbf{Linux, Kubernetes et automatisation cloud-native} avec \textbf{Infrastructure as Code, GitOps et DevSecOps}. Expérience dans le déploiement de plateformes \textbf{MLOps multi-tenant} et la conception d'agents IA pour l'automatisation des opérations. Auteur IEEE IECON 2026.

% ================================
% COMPÉTENCES TECHNIQUES
% ================================
\section{Compétences Techniques}

\cvitem{\textbf{Cloud \& Orchestration}}{
  \textbullet~Kubernetes · Docker · Proxmox · Azure · GCP
}

\cvitem{\textbf{IaC \& CI/CD}}{
  \textbullet~Terraform · Ansible · Helm · ArgoCD · GitLab CI · OpenTofu · Apache CloudStack
}

\cvitem{\textbf{MLOps \& IA Agentic}}{
  \textbullet~KServe · MLflow · LangChain · LangGraph · n8n
}

\cvitem{\textbf{DataOps \& Data Eng.}}{
  \textbullet~Airflow · Kafka · Apache Iceberg · Hive · Trino
}

\cvitem{\textbf{Observabilité \& Sécurité}}{
  \textbullet~Prometheus · Grafana · OpenTelemetry · Istio · Kyverno · Vault · Harbor · Trivy
}

\cvitem{\textbf{Langages \& Frameworks}}{
  \textbullet~Python · Bash · SQL · Django · Angular · Java
}

% ================================
% PROJETS PHARES
% ================================
\section{Projets Phares}

\cventry{2024}{KubeOlive MLOps Platform}{LAMBDA Lab}{}{}{
  Plateforme MLOps Charmed Kubeflow auto-hébergée sur Proxmox VE, automatisée avec Terraform et Ansible, avec stockage Longhorn, multi-tenant via GitLab OIDC, isolation MinIO par utilisateur et pipelines visuels Elyra.
  \newline
  \textbf{Stack :} Kubeflow · MicroK8s · Terraform · Ansible · Proxmox · Longhorn · MinIO · FastAPI
  \newline
  \textbf{GitHub :} \url{https://github.com/ahmad-bensalah/KubeOlive}
}

\cventry{2024}{DevAIPortal}{Projet Personnel}{}{}{
  Plateforme développeur interne conversationnelle qui transforme des demandes en langage naturel en environnements Kubernetes validés via GitOps, blueprints Helm et expiration automatique gérée par Kyverno.
  \newline
  \textbf{Stack :} Kubernetes · Helm · ArgoCD · Kyverno · LLM · TypeScript
  \newline
  \textbf{GitHub :} \url{https://github.com/ahmad-bensalah/DevAIPortal}
}

\cventry{2024}{TendOlive RFP Intelligence}{Challenge CSTAM 3.0}{}{}{
  Pipeline automatisé qui collecte les appels d'offres publics, vérifie chaque exigence avec un RAG sur la base de connaissances et génère le deck de proposition.
  \newline
  \textbf{Stack :} n8n · RAG · LLM · Firecrawl · JavaScript
  \newline
  \textbf{GitHub :} \url{https://github.com/ahmad-bensalah/TendOlive}
}

\cventry{2023}{SignBridge}{Projet Académique}{}{}{
  Plateforme basée sur l'IA assurant la traduction bidirectionnelle entre le dialecte tunisien et la langue des signes tunisienne.
  \newline
  \textbf{Stack :} BiLSTM · Mediapipe · TTS · STT · Python
  \newline
  \textbf{GitHub :} \url{https://github.com/ahmad-bensalah/SignBridge}
}

% ================================
% EXPÉRIENCE PROFESSIONNELLE
% ================================
\section{Expérience Professionnelle}

\cventry{Juil--Sep 2024}{DevOps \& MLOps Engineer}{LAMBDA Lab - FSS}{Sfax, Tunisie}{}{
  \begin{itemize}
    \item Déploiement d'une plateforme Kubeflow multi-tenant sur Proxmox avec IaC (Terraform, Ansible)
    \item Mise en place de GitOps avec ArgoCD et gestion du stockage distribué (Longhorn)
    \item Développement d'un portail self-service FastAPI pour l'onboarding des utilisateurs
    \item Formation de l'équipe aux workflows MLOps et pipelines Elyra
  \end{itemize}
}

\cventry{Mar--Mai 2024}{DataOps \& DevOps Intern}{InovareTech}{Sfax, Tunisie}{}{
  \begin{itemize}
    \item Automatisation du déploiement de clusters Hadoop avec Vagrant et Ansible
    \item Mise en place de pipelines CI/CD avec Jenkins, Docker et ArgoCD
    \item Configuration de l'observabilité avec Prometheus et Grafana
    \item Implémentation de politiques de sécurité Kyverno
  \end{itemize}
}

% ================================
% FORMATION
% ================================
\section{Formation}

\cventry{2022--2027}{Diplôme d'Ingénieur en Data Engineering}{Faculté des Sciences de Sfax}{Tunisie}{}{
  \textbf{Spécialisation :} DevOps, Cloud Computing, MLOps, Data Engineering
}

\cventry{2020--2022}{Classes Préparatoires MPSI/MP}{Lycée Pilote de Sfax}{Tunisie}{}{
  Mathématiques, Physique, Sciences de l'Ingénieur
}

% ================================
% PUBLICATIONS \& DISTINCTIONS
% ================================
\section{Publications \& Distinctions}

\cvitem{2026}{\textbf{IEEE IECON 2026 Author} - Article scientifique accepté sur l'IA agentic appliquée à l'aquaponie durable}

\cvitem{2024}{\textbf{CSTAM 3.0 Challenge} - Finaliste avec TendOlive (automatisation d'appels d'offres par IA)}

\cvitem{2023}{\textbf{IEEE TSYP13 Challenge} - Projet Aqura AI (plateforme d'aquaponie pilotée par la voix)}

% ================================
% LANGUES
% ================================
\section{Langues}

\cvitemwithcomment{Arabe}{Langue maternelle}{}
\cvitemwithcomment{Français}{Courant (C1)}{}
\cvitemwithcomment{Anglais}{Courant (C1)}{}

% ================================
% CERTIFICATIONS (Optional)
% ================================
% \section{Certifications}
% \cvitem{2024}{Certified Kubernetes Administrator (CKA) - En préparation}
% \cvitem{2024}{HashiCorp Certified: Terraform Associate - En préparation}

\end{document}
